\chapter{The Mirror That Changed Everything}
\label{stone:mirrors}

Children understand mirrors long before they understand mathematics. Stand in front of one. Raise your hand. Your reflection raises its hand. Walk toward it. It walks toward you.

A mirror changes everything. And yet somehow changes nothing. Distances are preserved. Lengths are preserved. Only one thing changes. The world is turned inside out.

If PhiBase really describes geometry, it should understand mirrors too.

\section*{A Builder's Question}

Imagine that we have completed a beautiful structure. Now place a perfectly flat mirror through its centre. Where does every beam go? Where does every corner go? Can the arithmetic describe the reflected structure just as exactly as it described the original one? That is our next test.

\section*{The Six Working Mirrors}

The dodecahedron contains many planes of symmetry. For our construction we begin with six of them. Each passes through the centre of the structure. Each is parallel to one family of dodecahedral faces. Those six mirrors are enough to reveal something extraordinary. Later we shall discover the complete family of fifteen. Today we need only these six.

\section*{The First Reflection}

Take one point. Reflect it. Now convert the reflected point back into a six-beam construction. Repeat the experiment. Again. Again.

At first everything seems ordinary. Then something begins happening often enough that it can no longer be ignored. The number \(5\) keeps appearing. Not once. Not twice. Almost every reflection seems to introduce a denominator of five.

That is too regular to be an accident. The geometry is trying to tell us something.

\section*{Where Did The Five Come From?}

The answer is hidden inside the mirror itself. Every mirror has a direction pointing straight away from its surface. Mathematicians call that direction the \textbf{normal}. For each of our six working mirrors, the normal has exactly the same PhiBase length,
\[
\varphi+2.
\]
Yesterday we learned that every PhiBase number has a conjugate. Today we use it again. The norm of \(\varphi+2\) turns out to be
\[
5.
\]
The denominator did not appear because of the arithmetic. It appeared because the geometry demanded it. The arithmetic merely reported what the geometry already knew.

\section*{A New Layer Of The Landscape}

At first a denominator of five feels disappointing. Have we left PhiBase? No. We have simply discovered a larger landscape. The reflected point is still exact. Nothing has been rounded. Nothing has been approximated. The point now belongs to
\[
\frac{1}{5}\mathbb{Z}[\varphi].
\]
The arithmetic has not broken. It has grown.

\section*{Can We Still Come Home?}

By now you already know the question we should ask. Can we reverse the construction? Sometimes the answer is immediate. Sometimes it is not. The arithmetic tells us exactly which reflected points belong to the original six-beam lattice. When the necessary combinations are divisible by ten, every beam count returns as a whole number. The construction comes home again. Otherwise the reflected point remains perfectly exact, but it lies between the ordinary lattice points. The geometry has become finer. The arithmetic has faithfully followed it.

\section*{A Familiar Friend}

Some points never move at all. Reflect them once. Reflect them a hundred times. They remain exactly where they began. Why? Because they already lie in the mirror. Those fixed points become one of our simplest ways to test the entire reflection machinery. Nature provides the answer before we ask the question.

\section*{Builder's Notebook}

Choose one point from a small construction. Reflect it. Then reflect the result. Did you return exactly to your starting point? If not, something is wrong. A mirror should always undo itself. Reflection is its own inverse. Just like the best ideas in mathematics.

\section*{Looking Ahead}

The mirrors have revealed something unexpected. PhiBase is larger than the original lattice. It can describe exact constructions that lie beyond the ordinary building points. One final question remains. Could a computer use this arithmetic just as naturally as we do? That is the last great construction before we return to the Builder's Way.
